\documentclass[10pt,a4paper]{amsart}
\usepackage[margin=19mm]{geometry}
\usepackage{amsmath,amssymb,mathtools}
\usepackage[scr=rsfs,cal=euler]{mathalfa}
\usepackage{xfrac}
\usepackage[unicode,hidelinks,bookmarks=false]{hyperref}
\newcommand{\ie}{i.e.\ }
\newcommand{\cf}{cf.\ }
\newcommand{\resp}{respectively\ }
\newcommand{\Subsection}{\subsection}
\providecommand{\nobreakdash}{\nobreak\hbox{-}}
\setlength{\emergencystretch}{2em}
\multlinegap0pt
\usepackage{bm}
\usepackage{bbm}
\usepackage{dsfont}
\usepackage{xspace}
\usepackage{cancel}
\usepackage{mathtools}
\usepackage{amsthm}
\makeatletter
\@ifundefined{theorem}{\newtheorem{theorem}{Theorem}}{}
\makeatother
\renewcommand{\geq}{\geqslant}
\renewcommand{\leq}{\leqslant}
\title[Extending the Affirmative Action Problem: mixing numbers and]{Extending the Affirmative Action Problem: mixing numbers and integrated colorings of graphs}
\author{Charles Burnette, Broden Caton, Olivia Coward, Julian Davis, Austin Teter}
\date{Source: arXiv:2506.07192v1 (CC BY 4.0)}
\begin{document}
\maketitle

\noindent\emph{Source and licence.} Extending the Affirmative Action Problem: mixing numbers and integrated colorings of graphs. Version arXiv:2506.07192v1 (CC BY 4.0), distributed under the Creative Commons Attribution 4.0 International licence, as recorded in the arXiv licence metadata for this exact version and shown on the abstract page https://arxiv.org/abs/2506.07192v1. Full rights notice: licenses/arxiv-2506.07192-CC-BY-4.0.txt.\par\smallskip

\noindent\textit{Editorial scope.} Counted unit: the Theorem whose source label is \texttt{probabilisticmethod} in the article (source0.tex lines 703--714, source label \texttt{probabilisticmethod}), delivered together with the article's own complete original proof of it (source0.tex lines 716--746, one \texttt{proof} environment of 31 lines). No mathematics is written, repaired or re-derived: statement and proof are the article's own text line for line. Required context retained uncounted: none. Every definition, display and label of those lines is retained unchanged. Omitted from this package and not redistributed: 1--702, 715--715, 747--889. Nothing inside a retained excerpt is omitted or altered: every retained body is delivered exactly as the article's lines are written, so no omission receipt of any kind accompanies this package. Citations: the retained excerpts carry no citation command at all, so no bibliography entry of the article is transcribed and no reference is redistributed. Reference adaptation: none was needed and none was made; every reference inside the delivered unit resolves inside it, so no unresolved reference or citation is produced. Printed slips kept literal: no printed word, symbol, hypothesis or number of the counted statement or of its proof was altered. Layout and class: the article's own class is \texttt{amsart} with packages including graphicx, color, float, tikz, amsmath, amssymb, amsthm, amscd, enumerate, amsfonts, fancyhdr, mathtools, framed, eucal; the wrapper is the lane's pinned \texttt{amsart} editorial wrapper at 10pt on A4, which restates the theorem-like environments the retained text needs and adds the article's own macro definitions that the retained text needs (\texttt{\textbackslash geq}, \texttt{\textbackslash leq}), with extra wrapper packages (bm, bbm, dsfont, xspace, cancel, mathtools). The delivered numbering is the unit's own. Assets: no retained span contains a figure environment, a graphics-inclusion command, a PDF-page inclusion, a table or a TikZ picture; no image file is copied into this package, the package declares no asset dependency and no image file is redistributed with it. Licence: version arXiv:2506.07192v1 of this article is distributed under the Creative Commons Attribution 4.0 International licence, as recorded in the arXiv licence metadata for this exact version and shown on the abstract page https://arxiv.org/abs/2506.07192v1. A full rights notice is delivered at licenses/arxiv-2506.07192-CC-BY-4.0.txt. Characters: every retained excerpt is pure ASCII, so the delivered engine, XeLaTeX with its default Latin Modern text font, reports no missing character.
\begin{theorem}
\label{probabilisticmethod}
Let $G$ be a simple graph with $\Delta(G) \geq 2.$ Then
\[\textup{ic}(G) \leq \frac{\sigma^2}{\sigma^2+(|V''|-\mu)^2}\cdot2^{|V'|},\]
where
\[\mu = \sum_{v \in V''} \sum_{k = \lceil\lambda_v-\frac{1}{2}\textup{deg}(v)\rceil}^{\lambda_v} \binom{\lambda_v}{k}2^{-\lambda_v}\]
and
\begin{align*}
\sigma^2 &= \mu + \sum_{\substack{\textup{distinct}\ v, w \in V'' \\ v,w\ \textup{adjacent}}}(\alpha_{0,1}(v,w) + \alpha_{1,1}(v,w)) 
\\ & \hspace{2em}+ \sum_{\substack{\textup{distinct}\ v, w \in V'' \\ v,w\ \textup{not adjacent}}} (\alpha_{0,0}(v,w) + \alpha_{1,0}(v,w)) - \mu^2.
\end{align*}
\end{theorem}

\begin{proof}
Consider a semi-random coloring $\mathcal{C}$ of $G$ obtained by coloring each vertex in $V'$ independently with the choice of black or white being equally likely, then assign any pendant vertices the opposite color of their neighbors. Observe that any vertex in $V\backslash V''$ is then automatically integrated. For each $v \in V'',$ let $I_v$ be the indicator random variable for the event that $v$ is integrated. If $X =\ \sum_{v \in V''} I_v,$ then by linearity of expectation along with the fact that for every $v \in V'',$ the vertex $v$ is guaranteed at least $\text{deg}(v) - \lambda_v$ neighbors of the opposite color,
\begin{equation}
\mathbb{E}X = \sum_{v \in V''} \mathbb{E}I_v = \sum_{v \in V''} \text{Pr}\!\left[\text{mix}(v) \geq \frac{1}{2}\text{deg}(v)\right] = \mu,
\end{equation}
and
\begin{align}
\mathbb{E}[X^2] &= \mu + \sum_{\text{distinct}\ v, w\in V''} 2\mathbb{E}[I_vI_w]
\\ &= \mu + \sum_{\substack{\text{distinct}\ v, w \in V'' \\ v,w\ \text{adjacent}}} 2\mathbb{E}[I_vI_w] + \sum_{\substack{\text{distinct}\ v, w \in V'' \\ v,w\ \text{not adjacent}}} 2\mathbb{E}[I_vI_w].
\end{align}
Now for each pair of distinct vertices $v, w \in V'',$ let $J_{v,w}$ be the indicator random variable for the event that $v$ and $w$ are different colors. It is easy to see that $J_{v,w}$ amounts to a fair Bernoulli trial, and
\begin{align}
\mathbb{E}[I_vI_w] &= \text{Pr}[I_vI_w = 1]
\\ &= \text{Pr}[I_vI_w = 1\ \text{and}\ J_{v,w}=0] + \text{Pr}[I_vI_w = 1\ \text{and}\ J_{v,w}=1].
\end{align}
If $v$ and $w$ are adjacent, then
\[\mathbb{E}[I_vI_w] = \frac{1}{2}\alpha_{0,1}(v,w) + \frac{1}{2}\alpha_{1,1}(v,w).\]
On the other hand, if $v$ and $w$ are not adjacent, then
\[\mathbb{E}[I_vI_w] = \frac{1}{2}\alpha_{0,0}(v,w) + \frac{1}{2}\alpha_{1,0}(v,w).\]
Piecing everything together, it follows that
\begin{align}
\mathbb{E}[X^2] &= \mu + \sum_{\substack{\textup{distinct}\ v, w \in V'' \\ v,w\ \textup{adjacent}}}(\alpha_{0,1}(v,w) + \alpha_{1,1}(v,w)) \nonumber
\\ & \hspace{2em}+ \sum_{\substack{\textup{distinct}\ v, w \in V'' \\ v,w\ \textup{not adjacent}}} (\alpha_{0,0}(v,w) + \alpha_{1,0}(v,w))
\end{align}
and so $\mathbb{V}X = \sigma^2.$ Therefore, by the Chebyshev-Cantelli inequality,
\begin{align}
\text{Pr}[\mathcal{C}\ \text{is integrated}] &= \text{Pr}[X = |V''|]
\\ &= \text{Pr}[X - \mu \geq |V''| - \mu] \leq \frac{\sigma^2}{\sigma^2+(|V''|-\mu)^2}.
\end{align}
Multiply the above inequality by the total number of semi-random colorings to complete the proof.
\end{proof}

\end{document}
